% META: {"id": "1SPE-SUITES-CO-009", "chapitre": "1SPE-SUITES", "type_objet": "corrige", "exercice_ref": "1SPE-SUITES-EX-009", "capacites_codes": ["C3"], "sources_inspiration": [], "status": "generated"}
% BEGIN-VERIFY
% from sympy import *
% n = symbols('n', integer=True, nonnegative=True)
% u = 6 * 2**n
% assert u.subs(n, 0) == 6
% assert u.subs(n, 1) == 12
% assert u.subs(n, 2) == 24
% assert u.subs(n, 3) == 48
% # Les deux premiers ecarts sont differents : la suite n'est pas arithmetique.
% assert u.subs(n, 1) - u.subs(n, 0) == 6
% assert u.subs(n, 2) - u.subs(n, 1) == 12
% # Quotient u_{n+1}/u_n
% quotient = u.subs(n, n+1) / u
% assert simplify(quotient - 2) == 0
% # Le rang 5 est le premier, et l'unique, rang auquel u_n vaut 192.
% assert u.subs(n, 5) == 192
% assert all(u.subs(n, k) != 192 for k in range(5))
% END-VERIFY
\begin{corrige}{1SPE-SUITES-EX-009}

\commentaireMarge{Calcul direct des premiers termes par substitution dans $u_n = 6 \times 2^n$.}

\textbf{Question 1.} On substitue $n = 0$, $1$, $2$, $3$ dans $u_n = 6 \times 2^n$ :

\[
u_0 = 6 \times 2^0 = 6 \times 1 = 6.
\]
\[
u_1 = 6 \times 2^1 = 6 \times 2 = 12.
\]
\[
u_2 = 6 \times 2^2 = 6 \times 4 = 24.
\]
\[
u_3 = 6 \times 2^3 = 6 \times 8 = 48.
\]

Donc $u_0 = 6$, $u_1 = 12$, $u_2 = 24$ et $u_3 = 48$.

\commentaireMarge{Une seule différence ne suffit pas à reconnaître une suite arithmétique.}

\textbf{Question 2.} Les deux premières différences valent
\[
u_1-u_0=12-6=6,
\qquad
u_2-u_1=24-12=12.
\]
Elles sont différentes : les différences ne sont pas constantes. L'affirmation
de Lina est donc fausse et la suite n'est pas arithmétique.

\commentaireMarge{Le quotient de deux termes consécutifs permet ici de reconnaître le modèle multiplicatif.}

\textbf{Question 3.} Puisque $u_n = 6 \times 2^n > 0$ pour tout $n \in \mathbb{N}$, le quotient $\dfrac{u_{n+1}}{u_n}$ est bien défini. On calcule :
\[
\frac{u_{n+1}}{u_n} = \frac{6 \times 2^{n+1}}{6 \times 2^n} = 2^{n+1-n} = 2.
\]

Ce quotient est constant, égal à $2$, pour tout $n \in \mathbb{N}$.

Donc la suite $(u_n)$ est une suite \textbf{géométrique} de raison $q = 2$.

\commentaireMarge{La recherche d'un rang exact revient ici à reconnaître une puissance de $2$.}

\textbf{Question 4.} On résout
\[
u_n=192
\quad\Longleftrightarrow\quad
6\times 2^n=192
\quad\Longleftrightarrow\quad
2^n=32=2^5.
\]
Ainsi, le rang cherché est $n=5$.

\end{corrige}
